% Options for packages loaded elsewhere
\PassOptionsToPackage{unicode}{hyperref}
\PassOptionsToPackage{hyphens}{url}
\documentclass[
  french,
]{article}
\usepackage{xcolor}
\usepackage{amsmath,amssymb}
\setcounter{secnumdepth}{-\maxdimen} % remove section numbering
\usepackage{iftex}
\ifPDFTeX
  \usepackage[T1]{fontenc}
  \usepackage[utf8]{inputenc}
  \usepackage{textcomp} % provide euro and other symbols
\else % if luatex or xetex
  \usepackage{unicode-math} % this also loads fontspec
  \defaultfontfeatures{Scale=MatchLowercase}
  \defaultfontfeatures[\rmfamily]{Ligatures=TeX,Scale=1}
\fi
\usepackage{lmodern}
\ifPDFTeX\else
  % xetex/luatex font selection
\fi
% Use upquote if available, for straight quotes in verbatim environments
\IfFileExists{upquote.sty}{\usepackage{upquote}}{}
\IfFileExists{microtype.sty}{% use microtype if available
  \usepackage[]{microtype}
  \UseMicrotypeSet[protrusion]{basicmath} % disable protrusion for tt fonts
}{}
\makeatletter
\@ifundefined{KOMAClassName}{% if non-KOMA class
  \IfFileExists{parskip.sty}{%
    \usepackage{parskip}
  }{% else
    \setlength{\parindent}{0pt}
    \setlength{\parskip}{6pt plus 2pt minus 1pt}}
}{% if KOMA class
  \KOMAoptions{parskip=half}}
\makeatother
\usepackage{color}
\usepackage{fancyvrb}
\newcommand{\VerbBar}{|}
\newcommand{\VERB}{\Verb[commandchars=\\\{\}]}
\DefineVerbatimEnvironment{Highlighting}{Verbatim}{commandchars=\\\{\}}
% Add ',fontsize=\small' for more characters per line
\newenvironment{Shaded}{}{}
\newcommand{\AlertTok}[1]{\textcolor[rgb]{1.00,0.00,0.00}{\textbf{#1}}}
\newcommand{\AnnotationTok}[1]{\textcolor[rgb]{0.38,0.63,0.69}{\textbf{\textit{#1}}}}
\newcommand{\AttributeTok}[1]{\textcolor[rgb]{0.49,0.56,0.16}{#1}}
\newcommand{\BaseNTok}[1]{\textcolor[rgb]{0.25,0.63,0.44}{#1}}
\newcommand{\BuiltInTok}[1]{\textcolor[rgb]{0.00,0.50,0.00}{#1}}
\newcommand{\CharTok}[1]{\textcolor[rgb]{0.25,0.44,0.63}{#1}}
\newcommand{\CommentTok}[1]{\textcolor[rgb]{0.38,0.63,0.69}{\textit{#1}}}
\newcommand{\CommentVarTok}[1]{\textcolor[rgb]{0.38,0.63,0.69}{\textbf{\textit{#1}}}}
\newcommand{\ConstantTok}[1]{\textcolor[rgb]{0.53,0.00,0.00}{#1}}
\newcommand{\ControlFlowTok}[1]{\textcolor[rgb]{0.00,0.44,0.13}{\textbf{#1}}}
\newcommand{\DataTypeTok}[1]{\textcolor[rgb]{0.56,0.13,0.00}{#1}}
\newcommand{\DecValTok}[1]{\textcolor[rgb]{0.25,0.63,0.44}{#1}}
\newcommand{\DocumentationTok}[1]{\textcolor[rgb]{0.73,0.13,0.13}{\textit{#1}}}
\newcommand{\ErrorTok}[1]{\textcolor[rgb]{1.00,0.00,0.00}{\textbf{#1}}}
\newcommand{\ExtensionTok}[1]{#1}
\newcommand{\FloatTok}[1]{\textcolor[rgb]{0.25,0.63,0.44}{#1}}
\newcommand{\FunctionTok}[1]{\textcolor[rgb]{0.02,0.16,0.49}{#1}}
\newcommand{\ImportTok}[1]{\textcolor[rgb]{0.00,0.50,0.00}{\textbf{#1}}}
\newcommand{\InformationTok}[1]{\textcolor[rgb]{0.38,0.63,0.69}{\textbf{\textit{#1}}}}
\newcommand{\KeywordTok}[1]{\textcolor[rgb]{0.00,0.44,0.13}{\textbf{#1}}}
\newcommand{\NormalTok}[1]{#1}
\newcommand{\OperatorTok}[1]{\textcolor[rgb]{0.40,0.40,0.40}{#1}}
\newcommand{\OtherTok}[1]{\textcolor[rgb]{0.00,0.44,0.13}{#1}}
\newcommand{\PreprocessorTok}[1]{\textcolor[rgb]{0.74,0.48,0.00}{#1}}
\newcommand{\RegionMarkerTok}[1]{#1}
\newcommand{\SpecialCharTok}[1]{\textcolor[rgb]{0.25,0.44,0.63}{#1}}
\newcommand{\SpecialStringTok}[1]{\textcolor[rgb]{0.73,0.40,0.53}{#1}}
\newcommand{\StringTok}[1]{\textcolor[rgb]{0.25,0.44,0.63}{#1}}
\newcommand{\VariableTok}[1]{\textcolor[rgb]{0.10,0.09,0.49}{#1}}
\newcommand{\VerbatimStringTok}[1]{\textcolor[rgb]{0.25,0.44,0.63}{#1}}
\newcommand{\WarningTok}[1]{\textcolor[rgb]{0.38,0.63,0.69}{\textbf{\textit{#1}}}}
\usepackage{longtable,booktabs,array}
\newcounter{none} % for unnumbered tables
\usepackage{calc} % for calculating minipage widths
% Correct order of tables after \paragraph or \subparagraph
\usepackage{etoolbox}
\makeatletter
\patchcmd\longtable{\par}{\if@noskipsec\mbox{}\fi\par}{}{}
\makeatother
% Allow footnotes in longtable head/foot
\IfFileExists{footnotehyper.sty}{\usepackage{footnotehyper}}{\usepackage{footnote}}
\makesavenoteenv{longtable}
\ifLuaTeX
\usepackage[bidi=basic,shorthands=off]{babel}
\else
\usepackage[bidi=default,shorthands=off]{babel}
\fi
\ifLuaTeX
  \usepackage{selnolig} % disable illegal ligatures
\fi
\setlength{\emergencystretch}{3em} % prevent overfull lines
\providecommand{\tightlist}{%
  \setlength{\itemsep}{0pt}\setlength{\parskip}{0pt}}
\usepackage{bookmark}
\IfFileExists{xurl.sty}{\usepackage{xurl}}{} % add URL line breaks if available
\urlstyle{same}
\hypersetup{
  pdftitle={Maximum de vraisemblance glouton et critère BIC pour la localisation de sources ponctuelles},
  pdflang={fr},
  hidelinks,
  pdfcreator={LaTeX via pandoc}}

\title{Maximum de vraisemblance glouton et critère BIC pour la
localisation de sources ponctuelles}
\usepackage{etoolbox}
\makeatletter
\providecommand{\subtitle}[1]{% add subtitle to \maketitle
  \apptocmd{\@title}{\par {\large #1 \par}}{}{}
}
\makeatother
\subtitle{Rapport sur la méthode greedy\_ml (2D) et sa généralisation en
3D (greedy\_ml\_nd)}
\author{}
\date{}

\begin{document}
\maketitle

\section{Maximum de vraisemblance + BIC : rapport
complet}\label{maximum-de-vraisemblance-bic-rapport-complet}

Branche \texttt{travail-sur-le-probleme-inverse}, dossier
\texttt{sparse\_sources/}. Ce rapport complète le compte rendu général
(\texttt{COMPTE\_RENDU.md}) : il est entièrement consacré à la méthode
qui donne les meilleurs résultats sur le scénario du projet, le
\textbf{maximum de vraisemblance glouton avec choix du nombre de sources
par le critère BIC} (abrégé « ML + BIC »). Les chiffres viennent des
fichiers CSV de \texttt{sparse\_sources/results/}, produits par les
scripts cités dans chaque section.

Les formules sont écrites en LaTeX (\texttt{\$...\$} et
\texttt{\$\$...\$\$}) ; le document se convertit en LaTeX ou en PDF avec
Pandoc (section 10).

\subsection{0. Résumé}\label{ruxe9sumuxe9}

\begin{enumerate}
\def\labelenumi{\arabic{enumi}.}
\tightlist
\item
  \textbf{La méthode.} On ajoute les sources une à une ; à chaque étape,
  toutes les positions et intensités sont réajustées \textbf{sans
  pénalité} en maximisant la vraisemblance gaussienne (bruit de chaque
  capteur pris en compte), et le nombre de sources est choisi par le
  critère \textbf{BIC} : une source n'est gardée que si elle fait
  baisser \(\chi^2\) de plus de \(3\log M\approx 12{,}5\) (en 2D, 64
  capteurs). Aucun paramètre à régler.
\item
  \textbf{C'est la meilleure méthode testée} sur le scénario du projet :
  elle bat Prony et le Frank-Wolfe dans toutes les cases (par exemple 80
  \% de réussite contre 55 \% et 40 \% pour 3 sources à
  \(\sigma = 10^{-4}\) ; 50 \% contre 5 \% et 0 \% pour 5 sources).
\item
  \textbf{Elle est optimale} : son erreur vaut 0,7 à 0,9 fois la borne
  de Cramér--Rao (CRB), en 2D comme en 3D, et elle réussit dans 90 à 100
  \% des tirages où la CRB le permet. On ne peut pas faire beaucoup
  mieux en changeant d'algorithme.
\item
  \textbf{La théorie du BIC est vérifiée} : une source est détectée dès
  que son rapport signal sur bruit dépasse environ 3 (probabilité
  mesurée égale à la prédiction à quelques points près), avec 2 \% de
  fausses alarmes sur du bruit seul. Surestimer le bruit est sans danger
  ; le sous-estimer d'un facteur 2 crée des sources parasites.
\item
  \textbf{Généralisation en 3D} (\texttt{ml3d.py}) : même méthode, 4
  paramètres par source. Avec des capteurs sur une sphère, 1 à 2 sources
  sont retrouvées à \(\sigma = 10^{-3}\) malgré la source centrale ; le
  BLASSO 3D du projet échoue partout. Un seul cercle de capteurs ne
  suffit pas en 3D (problème non identifiable).
\item
  \textbf{Ce qui limite, ce sont les données} : le bruit de la source
  centrale (CRB multipliée par 42 par rapport au même problème sans
  elle), le nombre de sources et leur proximité. Plus de capteurs aide
  lentement (CRB en \(1/\sqrt M\)) ; des capteurs en hauteur sur un
  cylindre n'aident pas.
\end{enumerate}

{\def\LTcaptype{none} % do not increment counter
\begin{longtable}[]{@{}
  >{\raggedright\arraybackslash}p{(\linewidth - 2\tabcolsep) * \real{0.5000}}
  >{\raggedright\arraybackslash}p{(\linewidth - 2\tabcolsep) * \real{0.5000}}@{}}
\toprule\noalign{}
\begin{minipage}[b]{\linewidth}\raggedright
situation (2D, source centrale du projet)
\end{minipage} & \begin{minipage}[b]{\linewidth}\raggedright
ce que ML + BIC retrouve
\end{minipage} \\
\midrule\noalign{}
\endhead
\bottomrule\noalign{}
\endlastfoot
\(\sigma = 10^{-4}\) & 1 à 3 sources de façon fiable (80 à 100 \%), 5
sources une fois sur deux \\
\(\sigma = 10^{-3}\) & 1 à 2 sources (75 à 100 \%), 3 sources dans 40 \%
des cas \\
\(\sigma = 10^{-2}\) & 1 source ; au-delà, la CRB dépasse la taille du
domaine \\
séparation minimale & 1 à \(\sigma = 10^{-5}\), 2 à
\(\sigma = 10^{-4}\), 4 à \(\sigma = 10^{-3}\) \\
intensité minimale (source isolée, \(\sigma = 10^{-3}\)) & environ 2
pour la détecter, 10 pour la localiser à 1 près \\
\end{longtable}
}

\subsection{1. Le problème}\label{le-probluxe8me}

\subsubsection{1.1 Modèle}\label{moduxe8le}

On cherche \(K\) sources ponctuelles de positions \(z_k\in D\) et
d'intensités \(I_k > 0\) à partir des mesures de \(M\) capteurs placés
en \(x_1,\dots,x_M\). Le noyau est celui du projet, la loi en \(1/r^2\)
:

\[
G(x_i, z) = \frac{1}{\lVert x_i - z\rVert^2 + h_i^2},
\qquad
y_i = \sum_{k=1}^{K} I_k\,G(x_i, z_k) + I_c\,G(x_i, c) .
\]

\begin{itemize}
\tightlist
\item
  En 2D, \(D = [0,10]\times[0,20]\) et les \(M = 64\) capteurs sont sur
  le cercle circonscrit à \(D\), dans le plan des sources (\(h_i = 0\)).
  La hauteur \(h_i\) sert seulement à l'étude des capteurs sur un
  cylindre (section 6.6).
\item
  En 3D, \(D = [0,10]^3\), \(h_i = 0\) et les capteurs sont hors d'un
  même plan (section 6.7).
\item
  \(I_c\,G(x_i, c)\) est la \textbf{source centrale} du projet :
  intensité \(I_c = 50\times 100 = 5000\), au centre \(c\) de \(D\).
\end{itemize}

On note
\(\varphi(z) = \big(G(x_1,z),\dots,G(x_M,z)\big)\in\mathbb{R}^M\) la
signature d'une source d'intensité 1 placée en \(z\), et
\(A(z, I) = \sum_k I_k\,\varphi(z_k)\).

\subsubsection{1.2 Bruit et source
centrale}\label{bruit-et-source-centrale}

Le bruit du projet est \textbf{relatif} : le capteur \(i\) mesure
\(y_i(1 + \sigma\varepsilon_i)\), avec
\(\varepsilon_i\sim\mathcal N(0,1)\) indépendants. La contribution de la
source centrale est connue et retirée :

\[
m_i = y_i(1 + \sigma\varepsilon_i) - I_c\,G(x_i, c)
    = A(z, I)_i + w_i,
\qquad w_i\sim\mathcal N\big(0,\,s_i^2\big),\quad s_i = \sigma\,y_i .
\]

Le signal de la source centrale disparaît, \textbf{pas son bruit} : en
2D, \(I_c\,G(x_i, c) = 5000/125 = 40\) pour chaque capteur, alors que le
signal utile des petites sources vaut environ 7 en moyenne quadratique.
L'écart type du bruit, \(s_i\approx 40\,\sigma\), est donc dominé par la
source centrale. C'est ce qui rend le problème difficile.

\subsubsection{1.3 Pourquoi pas le BLASSO}\label{pourquoi-pas-le-blasso}

Le BLASSO (étudié dans le compte rendu général) résout

\[
\min_{K,\,z,\,I\ge 0}\ \tfrac12\big\lVert m - A(z, I)\big\rVert^2
+ \lambda\sum_k I_k\,\lVert\varphi(z_k)\rVert^{\alpha} .
\]

Il a trois défauts sur ce problème :

\begin{enumerate}
\def\labelenumi{\arabic{enumi}.}
\tightlist
\item
  \textbf{un seul paramètre \(\lambda\) règle deux choses} : le nombre
  de sources (il faut \(\lambda\) assez grand pour rejeter le bruit) et
  le biais des positions et intensités (il faut \(\lambda\) petit pour
  ne pas les fausser) ;
\item
  il ne réussit que si la \textbf{condition de Duval--Peyré} (H2) est
  vraie : un certificat qui reste sous 1 hors des sources. Avec le noyau
  \(1/r^2\), elle est fausse dès que deux sources sont proches ou qu'une
  source est près d'un capteur ou d'un bord, ce qui arrive souvent à
  partir de 3 sources ;
\item
  il ignore la structure du bruit (la variance \(s_i^2\) est connue).
\end{enumerate}

Le maximum de vraisemblance + BIC \textbf{sépare} les deux rôles de
\(\lambda\) : positions et intensités sont ajustées sans pénalité, et le
nombre de sources est choisi par un critère qui ne dépend que de \(M\).

\subsection{2. La méthode}\label{la-muxe9thode}

\subsubsection{2.1 Vraisemblance et
blanchiment}\label{vraisemblance-et-blanchiment}

Le bruit étant gaussien, indépendant, de variance \(s_i^2\) connue, la
log-vraisemblance vaut, à une constante près,

\[
\log L(z, I) = -\frac12\sum_{i=1}^{M}\frac{\big(m_i - A(z, I)_i\big)^2}{s_i^2}
= -\frac12\,\chi^2(z, I),
\qquad
\chi^2(z, I) = \big\lVert m_w - A_w(z, I)\big\rVert^2 ,
\]

avec les mesures et le modèle \textbf{blanchis} : \((m_w)_i = m_i/s_i\)
et \(A_w(z, I)_i = A(z, I)_i/s_i\). Chaque capteur compte d'autant moins
que son bruit est grand. En pratique, \(s_i = \sigma\,y_i\) est calculé
sur les mesures brutes, source centrale comprise
(\texttt{relative\_noise\_std}).

Si le modèle est juste, \(\chi^2\) au vrai paramètre suit une loi
\(\chi^2_M\) : il vaut environ \(M\). C'est un contrôle simple de la
reconstruction (« \(\chi^2/M\approx 1\) »).

\subsubsection{2.2 À nombre de sources fixé : moindres carrés non
linéaires}\label{uxe0-nombre-de-sources-fixuxe9-moindres-carruxe9s-non-linuxe9aires}

Pour \(K\) fixé, l'estimateur du maximum de vraisemblance est

\[
(\hat z, \hat I)_K = \operatorname*{arg\,min}_{z_k\in D,\ I_k\ge 0}\ \chi^2(z, I) .
\]

C'est un problème de moindres carrés non linéaires en \(z\) (linéaire en
\(I\)). On le résout par \textbf{descente conjointe} des positions et
intensités (L-BFGS-B, contraintes \(z_k\in D\) et \(I_k\ge 0\), gradient
analytique : \(\partial_{I_k}\chi^2 = -2\,\varphi_w(z_k)^{\top} r_w\) et
\(\nabla_{z_k}\chi^2 = -2\,I_k\,\nabla\varphi_w(z_k)^{\top} r_w\), avec
\(r_w = m_w - A_w(z, I)\) le résidu blanchi). Il n'y a \textbf{aucune
pénalité} : pas de biais en \(O(\lambda)\).

Le problème n'est pas convexe ; d'où la construction gloutonne de la
section 2.4, qui fournit un bon point de départ.

\subsubsection{2.3 Choix du nombre de sources : le critère
BIC}\label{choix-du-nombre-de-sources-le-crituxe8re-bic}

\(\chi^2(\hat z, \hat I)_K\) diminue toujours quand \(K\) augmente : une
source de plus « explique » toujours un peu de bruit. On choisit donc
\(K\) en minimisant

\[
C(K) = \chi^2\big(\hat z, \hat I\big)_K + \kappa\, p\, K,
\qquad p = d + 1 \ \text{paramètres par source},
\]

avec \(d\) la dimension (\(p = 3\) en 2D : deux coordonnées et une
intensité ; \(p = 4\) en 3D) et

\begin{itemize}
\tightlist
\item
  \(\kappa = \log M\) : \textbf{BIC} (critère d'information bayésien,
  Schwarz 1978) ;
\item
  \(\kappa = 2\) : \textbf{AIC} (Akaike 1974).
\end{itemize}

\textbf{Pourquoi \(\log M\).} Le BIC approche la probabilité a
posteriori du modèle à \(K\) sources : par la méthode de Laplace,
\(\log p(m\mid K)\approx \log L(\hat\theta_K) - \tfrac{pK}{2}\log M\).
Chaque paramètre « coûte » \(\tfrac12\log M\) en log-vraisemblance, soit
\(\log M\) en \(\chi^2\). Le BIC est \textbf{consistant} : si le vrai
modèle fait partie des candidats, il est choisi avec une probabilité qui
tend vers 1 quand \(M\) grandit ; l'AIC, lui, surestime \(K\) avec une
probabilité non nulle.

\textbf{Le seuil, concrètement.} Une source supplémentaire est acceptée
si elle fait baisser \(\chi^2\) de plus de \(p\,\log M\), soit
\(3\log 64\approx 12{,}5\) en 2D avec 64 capteurs.

\begin{itemize}
\tightlist
\item
  Une \textbf{fausse source}, qui n'explique que du bruit, fait baisser
  \(\chi^2\) d'environ \(p = 3\) en moyenne (loi \(\chi^2_3\)). Elle est
  presque toujours rejetée. (La recherche du meilleur emplacement
  augmente un peu cette baisse, d'où un taux de fausses détections
  faible mais non nul.)
\item
  Une \textbf{vraie source} de rapport signal sur bruit
  \(\mathrm{SNR} = I\,\lVert\varphi_w(z)\rVert\) fait baisser \(\chi^2\)
  d'environ \(\mathrm{SNR}^2 + 3\) (loi \(\chi^2_3\) décentrée de
  paramètre \(\mathrm{SNR}^2\)). Elle est acceptée si
  \(\mathrm{SNR}^2 + 3 > 12{,}5\), soit \(\mathrm{SNR}\gtrsim 3\). La
  section 6.2 vérifie cette prédiction.
\end{itemize}

Il n'y a \textbf{aucun paramètre à régler} : \(\kappa\) ne dépend que de
\(M\), et le niveau de bruit \(\sigma\) est celui du capteur (la section
6.4 étudie ce qui se passe s'il est mal connu).

\subsubsection{2.4 Algorithme glouton}\label{algorithme-glouton}

\begin{Shaded}
\begin{Highlighting}[]
\NormalTok{entrées : mesures m, écarts types s, capteurs, domaine D}
\NormalTok{K = 0, meilleur = (C(0), aucune source)}
\NormalTok{répéter}
\NormalTok{    r\_w = résidu blanchi}
\NormalTok{    nouvelle position z* = argmax\_z \textless{} phi\_w(z) / ||phi\_w(z)||, r\_w \textgreater{}}
\NormalTok{                           (grille de D, puis L{-}BFGS{-}B depuis les meilleurs noeuds)}
\NormalTok{    intensités par moindres carrés positifs (NNLS)}
\NormalTok{    glissement : descente conjointe de toutes les positions et intensités}
\NormalTok{    suppression des intensités nulles, fusion des sources confondues}
\NormalTok{    si C(K) \textless{} C(meilleur) : meilleur = solution courante}
\NormalTok{    sinon, si C n\textquotesingle{}a pas baissé depuis 2 ajouts : arrêt}
\NormalTok{étape arrière : retirer une source tant que C baisse}
\NormalTok{sortie : meilleure solution (z, I, K)}
\end{Highlighting}
\end{Shaded}

\begin{itemize}
\tightlist
\item
  Le choix de la nouvelle position utilise la \textbf{corrélation
  normalisée} (noyau \(\varphi_w/\lVert\varphi_w\rVert\), c'est-à-dire
  \(\alpha = 1\) ; en 2D, il faut pour cela passer un \texttt{Setup}
  normalisé, \texttt{make\_setup(normalize=True)}, ce que font tous les
  scripts ; \texttt{greedy\_ml\_nd} normalise toujours) : sans
  normalisation, le maximum serait toujours collé à un capteur, où
  \(\lVert\varphi\rVert\) explose. C'est le seul endroit où la
  normalisation intervient.
\item
  La \textbf{patience de 2 ajouts} évite de s'arrêter sur un palier ;
  l'\textbf{étape arrière} corrige les ajouts inutiles (par exemple une
  source placée au milieu de deux vraies sources, puis rendue inutile
  par les suivantes).
\item
  Coût : \(K_{\max}\) ajouts, chacun avec une descente conjointe sur
  \(pK\) variables ; 0,1 à 0,5 s en 2D pour 1 à 5 sources, 1 à 5 s en
  3D.
\end{itemize}

\subsubsection{2.5 Borne de Cramér--Rao}\label{borne-de-cramuxe9rrao}

La borne de Cramér--Rao (CRB) donne l'erreur minimale de tout estimateur
sans biais. Avec \(\theta = (I_1,\dots,I_K, z_1,\dots,z_K)\) et la
matrice des dérivées du modèle

\[
\mathbf{D} = \Big[\varphi(z_1),\dots,\varphi(z_K)\;\Big\vert\;I_1\nabla\varphi(z_1),\dots,I_K\nabla\varphi(z_K)\Big]\in\mathbb{R}^{M\times pK},
\]

la matrice de Fisher vaut
\(F = \mathbf{D}^{\top}\operatorname{diag}(s_i^{-2})\,\mathbf{D}\) et,
pour tout estimateur sans biais,

\[
\mathbb{E}\,\lVert\hat z_k - z_k\rVert^2 \;\ge\; \sum_{j=1}^{d}\big(F^{-1}\big)_{z_k^{(j)} z_k^{(j)}} =: \mathrm{CRB}_k^2 .
\]

Dans les tableaux, « CRB » désigne \(\max_k\mathrm{CRB}_k\) (la source
la moins bien déterminée). Deux propriétés guident la lecture des
résultats :

\begin{itemize}
\tightlist
\item
  la CRB est \textbf{proportionnelle au bruit} (\(\sigma\), donc à
  \(I_c\)) et décroît comme \(1/\sqrt M\) avec le nombre de capteurs ;
\item
  le maximum de vraisemblance est \textbf{asymptotiquement efficace} :
  quand le bruit est petit, son erreur atteint la CRB. Si une méthode
  atteint la CRB, \textbf{on ne peut pas faire mieux en changeant
  d'algorithme} ; il faut changer les données.
\end{itemize}

Si \(F\) est singulière, le problème n'est pas identifiable (par exemple
un seul cercle de capteurs en 3D, section 6.7) : la CRB est infinie.

\subsubsection{2.6 Généralisation en dimension quelconque
(3D)}\label{guxe9nuxe9ralisation-en-dimension-quelconque-3d}

Rien dans les sections 2.1 à 2.5 ne dépend de la dimension.
\texttt{ml3d.py} reprend la méthode pour des sources dans
\(\mathbb{R}^d\) :

\begin{itemize}
\tightlist
\item
  noyau \texttt{InverseSquareKernel} du projet (il accepte déjà des
  capteurs en dimension quelconque) ;
\item
  \(p = d + 1\) paramètres par source dans le BIC ;
\item
  grille de recherche de \(D\) en \(d\) dimensions (pas = côté / 20 en
  3D, soit \(21^3 = 9261\) nœuds) ;
\item
  même descente conjointe, même étape arrière, même CRB.
\end{itemize}

\textbf{Identifiabilité en 3D.} Avec des sources en 3D et un seul cercle
de capteurs horizontal, une source au-dessus du plan du cercle et sa
symétrique en dessous donnent exactement les mêmes mesures : le problème
n'est pas identifiable, et \(F\) est singulière. Il faut des capteurs
non coplanaires (deux ou trois cercles à des hauteurs différentes, ou
une sphère). Le test \texttt{test\_single\_ring\_is\_not\_identifiable}
le vérifie.

En 2D, \texttt{greedy\_ml\_nd} donne exactement la même solution que
\texttt{greedy\_ml} et la même borne que \texttt{cramer\_rao} (test
\texttt{test\_matches\_2d\_implementation}).

\subsection{3. Implémentation}\label{impluxe9mentation}

{\def\LTcaptype{none} % do not increment counter
\begin{longtable}[]{@{}
  >{\raggedright\arraybackslash}p{(\linewidth - 2\tabcolsep) * \real{0.5000}}
  >{\raggedright\arraybackslash}p{(\linewidth - 2\tabcolsep) * \real{0.5000}}@{}}
\toprule\noalign{}
\begin{minipage}[b]{\linewidth}\raggedright
fichier
\end{minipage} & \begin{minipage}[b]{\linewidth}\raggedright
contenu
\end{minipage} \\
\midrule\noalign{}
\endhead
\bottomrule\noalign{}
\endlastfoot
\texttt{ml.py} & \texttt{greedy\_ml} (2D, utilise \texttt{Setup}),
\texttt{cramer\_rao}, \texttt{relative\_noise\_std} \\
\texttt{ml3d.py} & \texttt{greedy\_ml\_nd}, \texttt{cramer\_rao\_nd},
\texttt{match\_nd}, \texttt{success\_nd} (dimension quelconque) \\
\texttt{kernel.py} & \texttt{Setup} : blanchiment
(\texttt{with\_weights}), hauteur par capteur (\texttt{heights}),
\texttt{make\_cylinder\_setup} \\
\texttt{tester\_ml.py} & \textbf{script pour tester la méthode à la
main} sur un tirage (2D ou 3D) \\
\texttt{experiments\_projet.py} & comparaison avec Prony et Frank-Wolfe
sur le scénario du projet \\
\texttt{experiments\_ml.py} & études propres à ML + BIC : \(\sigma\) mal
connu, détection, résolution, nombre de capteurs \\
\texttt{experiments\_cylindre.py} & capteurs sur la paroi d'un cylindre,
\(K\) connu \\
\texttt{experiments\_ml3d.py} & sources en 3D \\
\texttt{tests/test\_ml.py}, \texttt{tests/test\_ml3d.py},
\texttt{tests/test\_cylinder.py} & tests pytest \\
\end{longtable}
}

\textbf{Utilisation en 2D :}

\begin{Shaded}
\begin{Highlighting}[]
\ImportTok{from}\NormalTok{ sparse\_sources }\ImportTok{import}\NormalTok{ greedy\_ml, cramer\_rao, relative\_noise\_std, make\_setup}

\NormalTok{S }\OperatorTok{=}\NormalTok{ make\_setup(h}\OperatorTok{=}\FloatTok{0.0}\NormalTok{, normalize}\OperatorTok{=}\VariableTok{True}\NormalTok{)       }\CommentTok{\# capteurs du projet}
\NormalTok{s }\OperatorTok{=}\NormalTok{ relative\_noise\_std(y\_brut, sigma)       }\CommentTok{\# s\_i = sigma * y\_i (source centrale comprise)}
\NormalTok{r }\OperatorTok{=}\NormalTok{ greedy\_ml(m, s, setup}\OperatorTok{=}\NormalTok{S)                }\CommentTok{\# m : mesures, source centrale retirée}
\NormalTok{r[}\StringTok{"z"}\NormalTok{], r[}\StringTok{"I"}\NormalTok{], r[}\StringTok{"K"}\NormalTok{], r[}\StringTok{"chi2"}\NormalTok{] }\OperatorTok{/}\NormalTok{ S.M     }\CommentTok{\# positions, intensités, K, chi2 / M}
\NormalTok{cramer\_rao(r[}\StringTok{"z"}\NormalTok{], r[}\StringTok{"I"}\NormalTok{], s, S)[}\StringTok{"pos\_std"}\NormalTok{] }\CommentTok{\# précision attendue de chaque source}
\end{Highlighting}
\end{Shaded}

\textbf{Utilisation en 3D :}

\begin{Shaded}
\begin{Highlighting}[]
\ImportTok{from}\NormalTok{ sparse\_sources.ml3d }\ImportTok{import}\NormalTok{ greedy\_ml\_nd, cramer\_rao\_nd}
\NormalTok{r }\OperatorTok{=}\NormalTok{ greedy\_ml\_nd(m, s, sensors, box)        }\CommentTok{\# sensors (M, 3), box : Box 3D}
\end{Highlighting}
\end{Shaded}

\textbf{Paramètres} (valeurs par défaut conseillées) :
\texttt{criterion="bic"} ; \texttt{k\_max=10} ; \texttt{patience=2} ;
\texttt{backward=True}. Pour imposer \(K\) :
\texttt{k\_max=K,\ criterion=0.0,\ backward=False}.

\subsection{4. Tests}\label{tests}

\subsubsection{4.1 Tests automatiques
(pytest)}\label{tests-automatiques-pytest}

\texttt{pytest\ sparse\_sources/tests} lance 86 tests (environ 70 s),
dont 24 consacrés à ML + BIC et à ses extensions. Tous passent.

\textbf{\texttt{tests/test\_ml.py}} (2D, 8 tests) :

{\def\LTcaptype{none} % do not increment counter
\begin{longtable}[]{@{}
  >{\raggedright\arraybackslash}p{(\linewidth - 2\tabcolsep) * \real{0.5000}}
  >{\raggedright\arraybackslash}p{(\linewidth - 2\tabcolsep) * \real{0.5000}}@{}}
\toprule\noalign{}
\begin{minipage}[b]{\linewidth}\raggedright
test
\end{minipage} & \begin{minipage}[b]{\linewidth}\raggedright
ce qui est vérifié
\end{minipage} \\
\midrule\noalign{}
\endhead
\bottomrule\noalign{}
\endlastfoot
\texttt{test\_whitened\_derivatives} & dérivées du noyau blanchi et
normalisé contre différences finies \\
\texttt{test\_weights\_api} & validation des poids, blanchiment,
invariance de \(\psi\) par des poids constants \\
\texttt{test\_greedy\_ml\_two\_sources} (\(\sigma = 10^{-4}\) et
\(10^{-3}\)) & 2 sources avec la source centrale : \(K\) exact, erreur
inférieure à 5 CRB, \(\chi^2\) compatible avec le bruit \\
\texttt{test\_greedy\_ml\_three\_sources\_low\_noise} & 3 sources
tirées, \(\sigma = 10^{-5}\) : \(K\) exact, positions précises \\
\texttt{test\_greedy\_ml\_options} & critères AIC et \(\kappa\)
numérique, candidats multiples, départ imposé (points hors de \(D\)
ignorés) \\
\texttt{test\_greedy\_ml\_pure\_noise} & bruit seul : au plus une
source \\
\texttt{test\_cramer\_rao\_scaling\_and\_efficiency} & CRB
proportionnelle au bruit ; sur 30 tirages, erreur rms du ML entre 0,6 et
1,5 fois la CRB (\textbf{efficacité}) \\
\end{longtable}
}

\textbf{\texttt{tests/test\_ml3d.py}} (dimension quelconque, 10 tests) :

{\def\LTcaptype{none} % do not increment counter
\begin{longtable}[]{@{}
  >{\raggedright\arraybackslash}p{(\linewidth - 2\tabcolsep) * \real{0.5000}}
  >{\raggedright\arraybackslash}p{(\linewidth - 2\tabcolsep) * \real{0.5000}}@{}}
\toprule\noalign{}
\begin{minipage}[b]{\linewidth}\raggedright
test
\end{minipage} & \begin{minipage}[b]{\linewidth}\raggedright
ce qui est vérifié
\end{minipage} \\
\midrule\noalign{}
\endhead
\bottomrule\noalign{}
\endlastfoot
\texttt{test\_box\_grid} & grille de \(D\) : bords compris, pas
respecté \\
\texttt{test\_recovers\_three\_sources\_3d} (sphère, deux cercles) & 3
sources en 3D, \(\sigma = 10^{-4}\) : \(K\) exact, erreur inférieure à 5
CRB, \(\chi^2\) compatible avec le bruit \\
\texttt{test\_with\_central\_source\_3d} & source centrale de 5000
retirée, 2 sources retrouvées \\
\texttt{test\_noise\_only\_gives\_no\_source} & bruit seul : au plus une
source \\
\texttt{test\_imposed\_k} & \(\kappa = 0\), \texttt{k\_max} \(= K\) :
exactement \(K\) sources \\
\texttt{test\_single\_ring\_is\_not\_identifiable} & un seul cercle en
3D : CRB infinie \\
\texttt{test\_matches\_2d\_implementation} & en 2D, même solution que
\texttt{greedy\_ml} (à \(10^{-4}\) près) et même CRB que
\texttt{cramer\_rao} \\
\texttt{test\_cramer\_rao\_scales\_with\_noise} & CRB exactement
proportionnelle au bruit \\
\texttt{test\_match\_nd} & appariement hongrois et erreurs \\
\end{longtable}
}

\textbf{\texttt{tests/test\_cylinder.py}} (capteurs à plusieurs
hauteurs, 6 tests) : noyau capteur par capteur
\(1/(\lVert x_i - z\rVert^2 + h_i^2)\), un seul étage identique à
l'ancien noyau, équivalence des hauteurs \(+h\) et \(-h\), dérivées
(\(\alpha = 0\) et \(1\)), conservation des hauteurs par
\texttt{with\_power}, \texttt{with\_normalization} et
\texttt{with\_weights}, validation des entrées.

\subsubsection{\texorpdfstring{4.2 Tester la méthode à la main :
\texttt{tester\_ml.py}}{4.2 Tester la méthode à la main : tester\_ml.py}}\label{tester-la-muxe9thode-uxe0-la-main-tester_ml.py}

\begin{Shaded}
\begin{Highlighting}[]
\ExtensionTok{python}\NormalTok{ sparse\_sources/tester\_ml.py                          }\CommentTok{\# 2D, 3 sources, sigma = 1e{-}3}
\ExtensionTok{python}\NormalTok{ sparse\_sources/tester\_ml.py }\AttributeTok{{-}{-}sources}\NormalTok{ 5 }\AttributeTok{{-}{-}sigma}\NormalTok{ 1e{-}4 }\AttributeTok{{-}{-}graine}\NormalTok{ 7}
\ExtensionTok{python}\NormalTok{ sparse\_sources/tester\_ml.py }\AttributeTok{{-}{-}positions} \StringTok{"2,4,30;7,15,70"} \AttributeTok{{-}{-}sigma}\NormalTok{ 1e{-}3}
\ExtensionTok{python}\NormalTok{ sparse\_sources/tester\_ml.py }\AttributeTok{{-}{-}sans{-}centrale} \AttributeTok{{-}{-}comparer}
\ExtensionTok{python}\NormalTok{ sparse\_sources/tester\_ml.py }\AttributeTok{{-}{-}dim}\NormalTok{ 3 }\AttributeTok{{-}{-}capteurs}\NormalTok{ sphere }\AttributeTok{{-}{-}sources}\NormalTok{ 3}
\end{Highlighting}
\end{Shaded}

Le script tire un jeu de sources (exactement comme
\texttt{scripts/reconstruct.py} en 2D), lance ML + BIC et affiche :

\begin{itemize}
\tightlist
\item
  les sources vraies avec leur CRB, puis les sources estimées ;
\item
  \(\chi^2/M\), qui doit être proche de 1 ;
\item
  l'évolution du critère \(C(K)\) à chaque ajout, ce qui montre où le
  BIC s'arrête ;
\item
  l'erreur et le verdict (réussite : bon \(K\) et erreur de position
  inférieure à \(\tfrac14\min(d_{\min}, 4)\), le critère de tout le
  dossier) ;
\item
  avec \texttt{-\/-comparer}, les résultats de Prony et du Frank-Wolfe
  sur les mêmes mesures.
\end{itemize}

Il enregistre une figure dans \texttt{outputs/tester\_ml/} (sources
vraies, estimées, cercles de rayon \(2\,\mathrm{CRB}\)). Exemple de
sortie (2D, 3 sources, \(\sigma = 10^{-3}\), graine 0) :

\begin{Shaded}
\begin{Highlighting}[]
\NormalTok{2D, 3 sources vraies, sigma = 0.001, source centrale : oui (5000, retirée), M = 64 capteurs}

\NormalTok{Sources vraies (et borne de Cramér{-}Rao sur la position) :}
\NormalTok{  position [6.37  5.396]  intensité   15.87  CRB 0.349}
\NormalTok{  position [0.41  0.331]  intensité   17.79  CRB 0.00415}
\NormalTok{  position [ 8.133 18.255]  intensité   66.34  CRB 0.0111}

\NormalTok{Maximum de vraisemblance + BIC : K = 3 (0.22 s), chi2 / M = 0.82}
\NormalTok{  position [0.406 0.322]  intensité   17.50}
\NormalTok{  position [ 8.14  18.269]  intensité   65.68}
\NormalTok{  position [5.813 5.721]  intensité   18.51}

\NormalTok{Évolution du critère : K=0: 735304.6, K=1: 238755.5, K=2: 997.1, K=3: 89.8,}
\NormalTok{                       K=4: 97.7, K=5: 102.9}

\NormalTok{Erreur de position max 0.645, erreur d\textquotesingle{}intensité relative max 0.167 {-}\textgreater{} RÉUSSITE}
\end{Highlighting}
\end{Shaded}

Le critère chute fortement jusqu'à \(K = 3\), puis remonte : chaque
source de plus coûte \(3\log 64\approx 12{,}5\) et ne fait baisser
\(\chi^2\) que de quelques unités. La source la plus faible, au centre
de \(D\) (près de la source centrale), est la moins bien localisée ;
c'est exactement ce qu'annonce sa CRB (0,35, contre 0,004 et 0,011 pour
les deux autres).

\subsection{5. Critère de réussite et
protocole}\label{crituxe8re-de-ruxe9ussite-et-protocole}

Sauf mention contraire : sources tirées comme dans \texttt{config/}
(positions uniformes, intensités uniformes sur le simplexe de somme
100), source centrale de 5000 ajoutée puis retirée, bruit relatif
\(\sigma\), mêmes mesures pour toutes les méthodes. Une reconstruction
est \textbf{réussie} si le nombre de sources est exact et si l'erreur de
position maximale (appariement hongrois) est inférieure à
\(\tfrac14\min(d_{\min}, 4)\), où \(d_{\min}\) est la plus petite
distance entre deux sources. Un tirage est \textbf{faisable} si
\(2\,\mathrm{CRB} < \tfrac14\min(d_{\min}, 4)\) : une méthode optimale a
alors de bonnes chances de réussir.

\subsection{6. Résultats}\label{ruxe9sultats}

\subsubsection{6.1 Comparaison avec Prony et le Frank-Wolfe (scénario du
projet)}\label{comparaison-avec-prony-et-le-frank-wolfe-scuxe9nario-du-projet}

\texttt{experiments\_projet.py}, 20 tirages par case, mêmes mesures pour
toutes les méthodes. Prony est la méthode de
\texttt{scripts/reconstruct.py} ; le Frank-Wolfe est le BLASSO pondéré
(\(\alpha = 1\), \(\lambda = 1{,}5\,\sigma\lVert y\rVert\)), la
meilleure variante du BLASSO trouvée dans le compte rendu général.

\textbf{Taux de réussite} (\texttt{projet\_methodes\_resume.csv}) :

{\def\LTcaptype{none} % do not increment counter
\begin{longtable}[]{@{}
  >{\raggedright\arraybackslash}p{(\linewidth - 12\tabcolsep) * \real{0.1429}}
  >{\raggedright\arraybackslash}p{(\linewidth - 12\tabcolsep) * \real{0.1429}}
  >{\raggedright\arraybackslash}p{(\linewidth - 12\tabcolsep) * \real{0.1429}}
  >{\raggedright\arraybackslash}p{(\linewidth - 12\tabcolsep) * \real{0.1429}}
  >{\raggedright\arraybackslash}p{(\linewidth - 12\tabcolsep) * \real{0.1429}}
  >{\raggedright\arraybackslash}p{(\linewidth - 12\tabcolsep) * \real{0.1429}}
  >{\raggedright\arraybackslash}p{(\linewidth - 12\tabcolsep) * \real{0.1429}}@{}}
\toprule\noalign{}
\begin{minipage}[b]{\linewidth}\raggedright
\(n\)
\end{minipage} & \begin{minipage}[b]{\linewidth}\raggedright
\(\sigma\)
\end{minipage} & \begin{minipage}[b]{\linewidth}\raggedright
Prony
\end{minipage} & \begin{minipage}[b]{\linewidth}\raggedright
Frank-Wolfe
\end{minipage} & \begin{minipage}[b]{\linewidth}\raggedright
\textbf{ML + BIC}
\end{minipage} & \begin{minipage}[b]{\linewidth}\raggedright
tirages faisables
\end{minipage} & \begin{minipage}[b]{\linewidth}\raggedright
ML + BIC sur les faisables
\end{minipage} \\
\midrule\noalign{}
\endhead
\bottomrule\noalign{}
\endlastfoot
1 & \(10^{-4}\) & 100 \% & 100 \% & \textbf{100 \%} & 100 \% & 100 \% \\
1 & \(10^{-3}\) & 15 \% & 100 \% & \textbf{100 \%} & 100 \% & 100 \% \\
1 & \(10^{-2}\) & 0 \% & 100 \% & \textbf{100 \%} & 85 \% & 100 \% \\
2 & \(10^{-4}\) & 85 \% & 90 \% & \textbf{95 \%} & 85 \% & 94 \% \\
2 & \(10^{-3}\) & 15 \% & 70 \% & \textbf{75 \%} & 80 \% & 94 \% \\
2 & \(3\times 10^{-3}\) & 0 \% & 40 \% & \textbf{75 \%} & 50 \% & 90
\% \\
2 & \(10^{-2}\) & 5 \% & 15 \% & \textbf{30 \%} & 20 \% & 100 \% \\
3 & \(10^{-4}\) & 55 \% & 40 \% & \textbf{80 \%} & 75 \% & 100 \% \\
3 & \(3\times 10^{-4}\) & 25 \% & 20 \% & \textbf{65 \%} & 35 \% & 100
\% \\
3 & \(10^{-3}\) & 10 \% & 5 \% & \textbf{40 \%} & 25 \% & 100 \% \\
5 & \(10^{-4}\) & 5 \% & 0 \% & \textbf{50 \%} & 20 \% & 100 \% \\
5 & \(3\times 10^{-4}\) & 5 \% & 0 \% & \textbf{15 \%} & 5 \% & 100
\% \\
\end{longtable}
}

\textbf{Erreur de position médiane et CRB} :

{\def\LTcaptype{none} % do not increment counter
\begin{longtable}[]{@{}llllll@{}}
\toprule\noalign{}
\(n\) & \(\sigma\) & Prony & Frank-Wolfe & ML + BIC & CRB médiane \\
\midrule\noalign{}
\endhead
\bottomrule\noalign{}
\endlastfoot
1 & \(10^{-3}\) & 0,059 & 0,023 & \textbf{0,022} & 0,025 \\
1 & \(10^{-2}\) & 0,84 & 0,23 & \textbf{0,23} & 0,25 \\
2 & \(10^{-4}\) & 0,014 & 0,035 & \textbf{0,014} & 0,016 \\
2 & \(10^{-3}\) & 0,32 & 0,29 & \textbf{0,13} & 0,16 \\
2 & \(3\times 10^{-3}\) & 1,6 & 0,66 & \textbf{0,36} & 0,49 \\
3 & \(10^{-4}\) & 0,14 & 0,72 & \textbf{0,15} & 0,22 \\
5 & \(10^{-4}\) & 0,88 & 1,6 & \textbf{0,60} & 1,04 \\
\end{longtable}
}

À retenir :

\begin{enumerate}
\def\labelenumi{\arabic{enumi}.}
\tightlist
\item
  ML + BIC est \textbf{la meilleure méthode dans toutes les cases}, ou à
  égalité avec la meilleure, et l'écart grandit avec le nombre de
  sources et le bruit (50 \% contre 5 \% et 0 \% pour 5 sources à
  \(\sigma = 10^{-4}\)).
\item
  \textbf{Son erreur atteint la CRB} : 0,7 à 0,9 fois la CRB médiane
  pour 1 et 2 sources. (La CRB du tableau porte sur la pire source,
  l'erreur médiane sur un tirage typique ; une erreur légèrement sous la
  CRB est donc normale.)
\item
  \textbf{Sur les tirages faisables, elle réussit dans 90 à 100 \% des
  cas.} Ses échecs sont presque tous des tirages où la CRB dépasse la
  tolérance : aucune méthode ne pourrait les réussir de façon fiable.
\end{enumerate}

\textbf{BIC contre AIC} (\texttt{projet\_critere.csv}) : le BIC est
meilleur dans toutes les cases testées (par exemple 80 \% contre 60 \%
pour 3 sources à \(\sigma = 10^{-4}\)). L'AIC ajoute des sources
parasites dans 10 à 20 \% des tirages ; les erreurs du BIC sont surtout
des sources \textbf{manquantes} à fort bruit, c'est-à-dire des sources
trop faibles pour être vues (section 6.2).

\textbf{Effet de la source centrale} (3 sources, \(\sigma = 10^{-3}\),
\texttt{projet\_centrale\_resume.csv}) :

{\def\LTcaptype{none} % do not increment counter
\begin{longtable}[]{@{}lll@{}}
\toprule\noalign{}
intensité centrale / 100 & ML + BIC & CRB médiane \\
\midrule\noalign{}
\endhead
\bottomrule\noalign{}
\endlastfoot
0 (pas de source centrale) & \textbf{100 \%} & 0,052 \\
5 & 70 \% & 0,30 \\
50 (projet) & 40 \% & 2,2 \\
500 & 0 \% & 21 \\
\end{longtable}
}

Sans source centrale, la méthode réussit toujours ; avec celle du
projet, la CRB est multipliée par environ 42. \textbf{La limite est le
bruit de la source centrale, pas l'algorithme.}

\subsubsection{6.2 Seuil de détection : la théorie du BIC est
vérifiée}\label{seuil-de-duxe9tection-la-thuxe9orie-du-bic-est-vuxe9rifiuxe9e}

\texttt{experiments\_ml.py}, expérience \texttt{detection} : une seule
source d'intensité \(I\) (0 à 20) à une position uniforme dans \(D\),
source centrale du projet, \(\sigma = 10^{-3}\), 40 tirages par
intensité (\texttt{ml\_detection.csv}, \texttt{ml\_detection.png}). La
section 2.3 prédit que la source est retenue avec la probabilité
\(P\big(\chi^2_3(\mathrm{SNR}^2) > 3\log M\big)\), avec
\(\mathrm{SNR} = I\,\lVert\varphi(z)/s\rVert\).

{\def\LTcaptype{none} % do not increment counter
\begin{longtable}[]{@{}
  >{\raggedright\arraybackslash}p{(\linewidth - 8\tabcolsep) * \real{0.2000}}
  >{\raggedright\arraybackslash}p{(\linewidth - 8\tabcolsep) * \real{0.2000}}
  >{\raggedright\arraybackslash}p{(\linewidth - 8\tabcolsep) * \real{0.2000}}
  >{\raggedright\arraybackslash}p{(\linewidth - 8\tabcolsep) * \real{0.2000}}
  >{\raggedright\arraybackslash}p{(\linewidth - 8\tabcolsep) * \real{0.2000}}@{}}
\toprule\noalign{}
\begin{minipage}[b]{\linewidth}\raggedright
SNR
\end{minipage} & \begin{minipage}[b]{\linewidth}\raggedright
tirages
\end{minipage} & \begin{minipage}[b]{\linewidth}\raggedright
détectée (\(K\ge 1\))
\end{minipage} & \begin{minipage}[b]{\linewidth}\raggedright
théorie
\end{minipage} & \begin{minipage}[b]{\linewidth}\raggedright
détectée et à moins de 1
\end{minipage} \\
\midrule\noalign{}
\endhead
\bottomrule\noalign{}
\endlastfoot
0 à 1 & 38 & 3 \% & 1 \% & 0 \% \\
1 à 2 & 34 & 6 \% & 6 \% & 0 \% \\
2 à 3 & 21 & 29 \% & 22 \% & 0 \% \\
3 à 4 & 17 & 65 \% & 63 \% & 12 \% \\
4 à 6 & 32 & 94 \% & 93 \% & 25 \% \\
6 à 10 & 36 & 100 \% & 100 \% & 53 \% \\
10 à 20 & 46 & 100 \% & 100 \% & 89 \% \\
plus de 20 & 96 & 100 \% & 100 \% & 100 \% \\
\end{longtable}
}

\begin{itemize}
\tightlist
\item
  \textbf{La probabilité de détection suit la théorie à quelques points
  près} : le seuil du BIC est bien \(\mathrm{SNR}\approx 3\) (50 \% de
  détection).
\item
  \textbf{Fausses alarmes} : sur 40 tirages de bruit seul (\(I = 0\)),
  une seule source parasite (2 \%). Le BIC ne crée presque jamais de
  source à partir du bruit.
\item
  \textbf{Détecter n'est pas localiser.} Entre \(\mathrm{SNR} = 3\) et
  10, la source est vue mais sa position est imprécise : la CRB vaut
  environ \(10/\mathrm{SNR}\) (4,2 pour \(I = 1\), 0,85 pour \(I = 5\)).
  Il faut \(\mathrm{SNR}\gtrsim 20\) pour une erreur inférieure à 1 dans
  presque tous les cas.
\item
  Dans le scénario du projet (\(\sigma = 10^{-3}\)), une source isolée
  d'intensité 2 est presque toujours détectée, mais il faut une
  intensité d'environ 10 pour la localiser à mieux que 1 près.
\end{itemize}

\subsubsection{6.3 Résolution de deux sources
proches}\label{ruxe9solution-de-deux-sources-proches}

Expérience \texttt{resolution} : deux sources d'intensité 50 à distance
\(d\), centre et orientation aléatoires dans \(D\), source centrale du
projet, 20 tirages (\texttt{ml\_resolution.csv},
\texttt{ml\_resolution.png}).

{\def\LTcaptype{none} % do not increment counter
\begin{longtable}[]{@{}
  >{\raggedright\arraybackslash}p{(\linewidth - 8\tabcolsep) * \real{0.2000}}
  >{\raggedright\arraybackslash}p{(\linewidth - 8\tabcolsep) * \real{0.2000}}
  >{\raggedright\arraybackslash}p{(\linewidth - 8\tabcolsep) * \real{0.2000}}
  >{\raggedright\arraybackslash}p{(\linewidth - 8\tabcolsep) * \real{0.2000}}
  >{\raggedright\arraybackslash}p{(\linewidth - 8\tabcolsep) * \real{0.2000}}@{}}
\toprule\noalign{}
\begin{minipage}[b]{\linewidth}\raggedright
\(d\)
\end{minipage} & \begin{minipage}[b]{\linewidth}\raggedright
\(\sigma = 10^{-5}\)
\end{minipage} & \begin{minipage}[b]{\linewidth}\raggedright
\(\sigma = 10^{-4}\)
\end{minipage} & \begin{minipage}[b]{\linewidth}\raggedright
\(\sigma = 10^{-3}\)
\end{minipage} & \begin{minipage}[b]{\linewidth}\raggedright
CRB médiane (\(\sigma = 10^{-3}\))
\end{minipage} \\
\midrule\noalign{}
\endhead
\bottomrule\noalign{}
\endlastfoot
0,5 & 80 \% & 20 \% & 10 \% & 7,8 \\
1 & \textbf{100 \%} & 70 \% & 15 \% & 2,1 \\
1,5 & 100 \% & 90 \% & 25 \% & 1,5 \\
2 & 100 \% & \textbf{100 \%} & 50 \% & 0,91 \\
3 & 100 \% & 100 \% & 90 \% & 0,47 \\
4 & 100 \% & 100 \% & \textbf{100 \%} & 0,36 \\
6 & 100 \% & 100 \% & 100 \% & 0,20 \\
\end{longtable}
}

\begin{itemize}
\tightlist
\item
  La \textbf{séparation minimale résolue} passe de 1
  (\(\sigma = 10^{-5}\)) à 2 (\(\sigma = 10^{-4}\)) puis 4
  (\(\sigma = 10^{-3}\)). La CRB d'une paire croît très vite quand \(d\)
  diminue (d'un facteur 2 à 4 chaque fois que \(d\) est divisé par 2) :
  deux sources proches ont presque la même signature.
\item
  Quand elle échoue, la méthode trouve \textbf{une seule source} au lieu
  de deux (85 \% des tirages à \(d = 0{,}5\), \(\sigma = 10^{-3}\)) : le
  BIC juge que la deuxième source n'apporte pas assez. C'est le
  comportement attendu, et il est préférable à l'invention de sources.
\item
  L'erreur médiane reste \textbf{sous la CRB} de la pire source à toutes
  les distances (figure de droite).
\end{itemize}

\subsubsection{6.4 Niveau de bruit mal
connu}\label{niveau-de-bruit-mal-connu}

La méthode a besoin des écarts types \(s_i = \sigma y_i\). Expérience
\texttt{sigma\_faux} : la méthode reçoit \(c\times\sigma\) au lieu de
\(\sigma\) (\texttt{ml\_sigma\_faux.csv}, \texttt{ml\_sigma\_faux.png},
20 tirages).

{\def\LTcaptype{none} % do not increment counter
\begin{longtable}[]{@{}
  >{\raggedright\arraybackslash}p{(\linewidth - 10\tabcolsep) * \real{0.1667}}
  >{\raggedright\arraybackslash}p{(\linewidth - 10\tabcolsep) * \real{0.1667}}
  >{\raggedright\arraybackslash}p{(\linewidth - 10\tabcolsep) * \real{0.1667}}
  >{\raggedright\arraybackslash}p{(\linewidth - 10\tabcolsep) * \real{0.1667}}
  >{\raggedright\arraybackslash}p{(\linewidth - 10\tabcolsep) * \real{0.1667}}
  >{\raggedright\arraybackslash}p{(\linewidth - 10\tabcolsep) * \real{0.1667}}@{}}
\toprule\noalign{}
\begin{minipage}[b]{\linewidth}\raggedright
cas
\end{minipage} & \begin{minipage}[b]{\linewidth}\raggedright
\(c = 0{,}5\)
\end{minipage} & \begin{minipage}[b]{\linewidth}\raggedright
\(c = 0{,}8\)
\end{minipage} & \begin{minipage}[b]{\linewidth}\raggedright
\(c = 1\)
\end{minipage} & \begin{minipage}[b]{\linewidth}\raggedright
\(c = 1{,}25\)
\end{minipage} & \begin{minipage}[b]{\linewidth}\raggedright
\(c = 2\)
\end{minipage} \\
\midrule\noalign{}
\endhead
\bottomrule\noalign{}
\endlastfoot
2 sources, \(\sigma = 10^{-3}\) & 5 \% (\(K\) = 3,0) & 80 \% & 75 \% &
80 \% & 80 \% (\(K\) = 1,8) \\
3 sources, \(\sigma = 10^{-4}\) & 25 \% (\(K\) = 4,1) & 70 \% & 80 \% &
75 \% & 75 \% (\(K\) = 2,9) \\
3 sources, \(\sigma = 10^{-3}\) & 10 \% (\(K\) = 3,8) & 35 \% & 40 \% &
40 \% & 35 \% (\(K\) = 2,4) \\
\end{longtable}
}

(entre parenthèses : nombre moyen de sources estimé)

\begin{itemize}
\tightlist
\item
  \textbf{Surestimer le bruit est sans danger} : jusqu'à un facteur 2,
  la réussite ne bouge presque pas. Le seuil effectif devient plus
  exigeant (\(\chi^2\) divisé par \(c^2\)) : on perd seulement les
  sources les plus faibles.
\item
  \textbf{Sous-estimer le bruit est dangereux} : avec \(c = 0{,}5\),
  \(\chi^2\) est multiplié par 4, chaque petite structure du bruit
  semble significative, et la méthode ajoute en moyenne une source
  parasite.
\item
  \textbf{Conseil} : en cas de doute, prendre \(\sigma\) un peu trop
  grand. Le contrôle \(\chi^2/M\) permet de vérifier le choix : avec le
  bon \(\sigma\), il vaut environ \((M - pK)/M\), soit 0,86 pour 3
  sources et 64 capteurs (0,78 à 0,85 mesuré). Une valeur nettement
  supérieure à 1 indique un \(\sigma\) trop petit ou une source
  manquante.
\end{itemize}

\subsubsection{6.5 Nombre de capteurs}\label{nombre-de-capteurs}

Expérience \texttt{capteurs} : même scénario, \(M\) capteurs sur le même
cercle, 20 tirages (\texttt{ml\_capteurs.csv},
\texttt{ml\_capteurs.png}).

{\def\LTcaptype{none} % do not increment counter
\begin{longtable}[]{@{}
  >{\raggedright\arraybackslash}p{(\linewidth - 8\tabcolsep) * \real{0.2000}}
  >{\raggedright\arraybackslash}p{(\linewidth - 8\tabcolsep) * \real{0.2000}}
  >{\raggedright\arraybackslash}p{(\linewidth - 8\tabcolsep) * \real{0.2000}}
  >{\raggedright\arraybackslash}p{(\linewidth - 8\tabcolsep) * \real{0.2000}}
  >{\raggedright\arraybackslash}p{(\linewidth - 8\tabcolsep) * \real{0.2000}}@{}}
\toprule\noalign{}
\begin{minipage}[b]{\linewidth}\raggedright
\(M\)
\end{minipage} & \begin{minipage}[b]{\linewidth}\raggedright
3 sources, \(\sigma = 10^{-3}\)
\end{minipage} & \begin{minipage}[b]{\linewidth}\raggedright
5 sources, \(\sigma = 10^{-4}\)
\end{minipage} & \begin{minipage}[b]{\linewidth}\raggedright
3 sources, \(\sigma = 10^{-2}\)
\end{minipage} & \begin{minipage}[b]{\linewidth}\raggedright
temps (5 sources)
\end{minipage} \\
\midrule\noalign{}
\endhead
\bottomrule\noalign{}
\endlastfoot
32 & 40 \% (CRB 3,7) & 35 \% (CRB 1,7) & 0 \% (CRB 36) & 0,6 s \\
64 (projet) & 40 \% (CRB 2,2) & 50 \% (CRB 1,0) & 0 \% (CRB 22) & 0,5
s \\
256 & 60 \% (CRB 1,1) & 65 \% (CRB 0,52) & 10 \% (CRB 11) & 1,0 s \\
1024 & \textbf{70 \%} (CRB 0,55) & \textbf{75 \%} (CRB 0,26) & 15 \%
(CRB 5,5) & 4,5 s \\
\end{longtable}
}

\begin{itemize}
\tightlist
\item
  La CRB est \textbf{divisée par 2 quand \(M\) est multiplié par 4} :
  elle décroît comme \(1/\sqrt M\), car les bruits des capteurs sont
  indépendants et se moyennent. ML + BIC en profite directement.
\item
  Le gain est réel mais lent : pour diviser l'erreur par 10, il faut 100
  fois plus de capteurs. À fort bruit (\(\sigma = 10^{-2}\)), même 1024
  capteurs ne suffisent pas.
\item
  Ce gain suppose des bruits \textbf{indépendants} d'un capteur à
  l'autre. Une part de bruit commune à tous les capteurs ne diminuerait
  pas.
\end{itemize}

\subsubsection{\texorpdfstring{6.6 Capteurs sur la paroi d'un cylindre
(\(K\)
connu)}{6.6 Capteurs sur la paroi d'un cylindre (K connu)}}\label{capteurs-sur-la-paroi-dun-cylindre-k-connu}

\texttt{experiments\_cylindre.py} (\texttt{cylindre.csv},
\texttt{cylindre\_resume.csv}, \texttt{cylindre.png}). Question posée :
avec des capteurs répartis sur la paroi d'un cylindre (plusieurs étages
au-dessus du plan des sources), le Frank-Wolfe marche-t-il pour beaucoup
de sources et beaucoup de bruit ?

\textbf{Protocole.} Même noyau \(1/(\lVert x_i - z\rVert^2 + h_i^2)\),
avec une hauteur \(h_i\) par capteur (\texttt{make\_cylinder\_setup}).
Le cylindre a pour axe la verticale du centre de \(D\) et pour rayon
celui du cercle du projet. Même nombre total de capteurs pour toutes les
géométries, \(M = 60\) :

{\def\LTcaptype{none} % do not increment counter
\begin{longtable}[]{@{}
  >{\raggedright\arraybackslash}p{(\linewidth - 2\tabcolsep) * \real{0.5000}}
  >{\raggedright\arraybackslash}p{(\linewidth - 2\tabcolsep) * \real{0.5000}}@{}}
\toprule\noalign{}
\begin{minipage}[b]{\linewidth}\raggedright
géométrie
\end{minipage} & \begin{minipage}[b]{\linewidth}\raggedright
étages (hauteurs \(h\))
\end{minipage} \\
\midrule\noalign{}
\endhead
\bottomrule\noalign{}
\endlastfoot
cercle & 1 étage, \(h = 0\) (comme le projet) \\
cylindre bas & \(h = 0\) ; 2,5 ; 5 (20 capteurs par étage) \\
cylindre haut & \(h = 0\) ; 2,5 ; 5 ; 7,5 ; 10 (12 par étage) \\
cylindre sans plan & \(h = 2{,}5\) ; 5 ; 7,5 (aucun capteur dans le
plan) \\
\end{longtable}
}

Comme les sources sont dans le plan, un capteur à \(-h\) mesure la même
chose qu'à \(+h\) : un cylindre symétrique de part et d'autre du plan
revient à ces étages avec deux fois plus de capteurs. Les étages sont
décalés en angle.

Sources : \(K\in\{3, 5, 8\}\), \textbf{\(K\) connu à l'avance},
positions uniformes mais \textbf{espacées} (distance minimale 3),
intensités uniformes sur le simplexe, source centrale de 5000, bruit
relatif, 15 tirages par case. Les mêmes sources sont utilisées pour
toutes les géométries.

\begin{itemize}
\tightlist
\item
  \textbf{Frank-Wolfe à \(\lambda\) optimal pour \(K\)} : BLASSO pondéré
  (\(\alpha = 1\)), 14 valeurs de \(\lambda\) de
  \(0{,}5\,\lambda_{\max}\) à \(3\times 10^{-4}\,\lambda_{\max}\) ;
  parmi celles qui donnent exactement \(K\) sources, on garde celle de
  plus petit résidu (après débiaisage). Si aucune ne donne \(K\) sources
  (29 \% des cas), on prend la première qui en donne plus et on garde
  les \(K\) plus intenses.
\item
  \textbf{ML à \(K\) imposé} :
  \texttt{greedy\_ml(k\_max=K,\ criterion=0,\ backward=False)}.
\end{itemize}

\textbf{Taux de réussite (Frank-Wolfe / ML), condition de Duval--Peyré
vraie, CRB médiane} :

{\def\LTcaptype{none} % do not increment counter
\begin{longtable}[]{@{}
  >{\raggedright\arraybackslash}p{(\linewidth - 10\tabcolsep) * \real{0.1667}}
  >{\raggedright\arraybackslash}p{(\linewidth - 10\tabcolsep) * \real{0.1667}}
  >{\raggedright\arraybackslash}p{(\linewidth - 10\tabcolsep) * \real{0.1667}}
  >{\raggedright\arraybackslash}p{(\linewidth - 10\tabcolsep) * \real{0.1667}}
  >{\raggedright\arraybackslash}p{(\linewidth - 10\tabcolsep) * \real{0.1667}}
  >{\raggedright\arraybackslash}p{(\linewidth - 10\tabcolsep) * \real{0.1667}}@{}}
\toprule\noalign{}
\begin{minipage}[b]{\linewidth}\raggedright
\(K\)
\end{minipage} & \begin{minipage}[b]{\linewidth}\raggedright
\(\sigma\)
\end{minipage} & \begin{minipage}[b]{\linewidth}\raggedright
cercle
\end{minipage} & \begin{minipage}[b]{\linewidth}\raggedright
cylindre bas
\end{minipage} & \begin{minipage}[b]{\linewidth}\raggedright
cylindre haut
\end{minipage} & \begin{minipage}[b]{\linewidth}\raggedright
cylindre sans plan
\end{minipage} \\
\midrule\noalign{}
\endhead
\bottomrule\noalign{}
\endlastfoot
3 & \(10^{-4}\) & \textbf{67 \%} / 87 \% & 53 \% / 87 \% & 53 \% /
\textbf{93 \%} & 53 \% / 87 \% \\
3 & \(10^{-3}\) & 20 \% / \textbf{53 \%} & 20 \% / 33 \% & 20 \% / 27 \%
& 7 \% / 27 \% \\
3 & \(10^{-2}\) & 7 \% / 7 \% & 7 \% / 0 \% & 0 \% / 0 \% & 7 \% / 0
\% \\
5 & \(10^{-4}\) & 0 \% / \textbf{33 \%} & 7 \% / 20 \% & 0 \% / 7 \% & 0
\% / 27 \% \\
5 & \(\ge 10^{-3}\) & 0 \% / 0 \% & 0 \% / 0 \% & 0 \% / 0 \% & 0 \% / 0
\% \\
8 & tous & 0 \% / 0 \% & 0 \% / 0 \% & 0 \% / 0 \% & 0 \% / 0 \% \\
condition vraie, \(K\) = 3 / 5 / 8 & & 47 / 0 / 0 \% & 47 / 7 / 0 \% &
53 / 7 / 0 \% & 47 / 7 / 0 \% \\
CRB médiane, \(K\) = 3 / 5 / 8 (\(\sigma = 10^{-4}\)) & & \textbf{0,16}
/ 4,5 / 19 & 0,20 / \textbf{2,5} / \textbf{6,8} & 0,29 / 3,6 / 11 & 0,28
/ 3,7 / 17 \\
\end{longtable}
}

\textbf{Réponse : non.} Le cylindre n'aide pas le Frank-Wolfe, ni pour
beaucoup de sources, ni pour beaucoup de bruit.

\begin{enumerate}
\def\labelenumi{\arabic{enumi}.}
\tightlist
\item
  \textbf{Le Frank-Wolfe n'est jamais meilleur sur un cylindre que sur
  le cercle}, même avec le \(\lambda\) choisi en connaissant \(K\) et
  des sources espacées. Il plafonne à 67 \% pour 3 sources à très faible
  bruit et tombe à 0 \% dès 5 sources, quelle que soit la géométrie.
\item
  \textbf{La condition de Duval--Peyré reste fausse} à partir de 5
  sources (0 à 7 \% des tirages). Les étages en hauteur améliorent un
  peu la marge (marge relative médiane pour 3 sources : \(-40\) pour le
  cercle, \(-17\) pour le cylindre bas, \(-9\) sans capteur dans le
  plan, \(+1{,}0\) pour le cylindre haut), car le signal ne peut plus
  exploser près d'un capteur. Mais ce gain ne se traduit pas en
  réussite.
\item
  \textbf{Monter les capteurs dégrade l'information à petit nombre de
  sources} : un capteur à la hauteur \(h\) reçoit moins de signal, et le
  reçoit plus « flou » (toutes les sources sont à peu près à la distance
  \(h\)). Pour 3 sources, la CRB passe de 0,16 (cercle) à 0,20, 0,28 et
  0,29. À 5 ou 8 sources, la diversité des points de vue aide un peu
  (CRB divisée par 2 à 3 avec le cylindre bas), mais la CRB reste très
  au-dessus de la tolérance : le problème reste insoluble pour toutes
  les méthodes.
\item
  \textbf{ML est au moins aussi bon que le Frank-Wolfe dans toutes les
  cases} (87 à 93 \% contre 53 à 67 \% pour 3 sources à
  \(\sigma = 10^{-4}\)).
\end{enumerate}

Ce qui limite, c'est le bruit de la source centrale et le nombre de
sources, pas la position des capteurs en hauteur. Pour un vrai gain, il
faudrait des capteurs \textbf{plus proches des sources} (à l'intérieur
de \(D\)) ou un bruit plus faible, pas des capteurs plus loin du plan.

\subsubsection{6.7 Généralisation en
3D}\label{guxe9nuxe9ralisation-en-3d}

\texttt{experiments\_ml3d.py} (\texttt{ml3d.csv},
\texttt{ml3d\_resume.csv}, \texttt{ml3d.png}) : sources dans le cube
\([0,10]^3\), espacées d'au moins 3, intensités uniformes sur le
simplexe (somme 100), source centrale de 5000 au centre du cube, bruit
relatif \(\sigma\), \(M = 160\) capteurs à 1,5 fois le rayon de la
sphère circonscrite (géométrie de \texttt{scripts/duval\_peyre\_3d.py}),
8 tirages par case (un tirage = 12,5 \%). Comparaison avec le BLASSO 3D
du projet (\texttt{particle\_sources.duval\_peyre.blasso}, noyau brut,
\(\lambda = 4\lVert
w\rVert\,\lVert\Phi\rVert\) et seuil de 0,5 \% comme dans
\texttt{scripts/duval\_peyre\_3d.py}).

\textbf{Taux de réussite de ML + BIC} (le BLASSO du projet est à 0 \%
dans toutes les cases sauf une source à \(\sigma = 10^{-3}\), 12 \%) :

{\def\LTcaptype{none} % do not increment counter
\begin{longtable}[]{@{}
  >{\raggedright\arraybackslash}p{(\linewidth - 10\tabcolsep) * \real{0.1667}}
  >{\raggedright\arraybackslash}p{(\linewidth - 10\tabcolsep) * \real{0.1667}}
  >{\raggedright\arraybackslash}p{(\linewidth - 10\tabcolsep) * \real{0.1667}}
  >{\raggedright\arraybackslash}p{(\linewidth - 10\tabcolsep) * \real{0.1667}}
  >{\raggedright\arraybackslash}p{(\linewidth - 10\tabcolsep) * \real{0.1667}}
  >{\raggedright\arraybackslash}p{(\linewidth - 10\tabcolsep) * \real{0.1667}}@{}}
\toprule\noalign{}
\begin{minipage}[b]{\linewidth}\raggedright
\(K\)
\end{minipage} & \begin{minipage}[b]{\linewidth}\raggedright
\(\sigma\)
\end{minipage} & \begin{minipage}[b]{\linewidth}\raggedright
deux cercles
\end{minipage} & \begin{minipage}[b]{\linewidth}\raggedright
trois cercles
\end{minipage} & \begin{minipage}[b]{\linewidth}\raggedright
sphère
\end{minipage} & \begin{minipage}[b]{\linewidth}\raggedright
CRB médiane (sphère)
\end{minipage} \\
\midrule\noalign{}
\endhead
\bottomrule\noalign{}
\endlastfoot
1 & \(10^{-4}\) & 100 \% & 100 \% & \textbf{100 \%} & 0,006 \\
1 & \(10^{-3}\) & 100 \% & 100 \% & \textbf{100 \%} & 0,059 \\
1 & \(10^{-2}\) & 88 \% & 88 \% & \textbf{100 \%} & 0,59 \\
2 & \(10^{-4}\) & 100 \% & 100 \% & \textbf{100 \%} & 0,049 \\
2 & \(10^{-3}\) & 62 \% & 62 \% & \textbf{88 \%} & 0,49 \\
2 & \(10^{-2}\) & 0 \% & 0 \% & 0 \% & 4,9 \\
3 & \(10^{-4}\) & 62 \% & 75 \% & \textbf{75 \%} & 0,23 \\
3 & \(10^{-3}\) & 12 \% & 25 \% & 25 \% & 2,3 \\
5 & \(10^{-4}\) & 0 \% & 12 \% & \textbf{50 \%} & 1,2 \\
5 & \(10^{-3}\) & 0 \% & 0 \% & 0 \% & 12 \\
\end{longtable}
}

\textbf{Erreur médiane et CRB} (une source) : 0,044 contre 0,059
(sphère, \(\sigma = 10^{-3}\)), 0,067 contre 0,085 (trois cercles),
0,074 contre 0,081 (deux cercles). L'erreur vaut 0,75 à 0,9 fois la CRB,
comme en 2D : \textbf{la méthode reste efficace en 3D.}

À retenir :

\begin{enumerate}
\def\labelenumi{\arabic{enumi}.}
\tightlist
\item
  \textbf{La méthode se généralise sans modification} : mêmes étapes,
  \(p = 4\) paramètres par source dans le BIC, mêmes garanties (erreur
  au niveau de la CRB, presque pas de fausses sources).
\item
  \textbf{La sphère est la meilleure disposition} : à nombre de capteurs
  égal, sa CRB est 1,4 fois (une source) à plus de 20 fois (cinq
  sources) plus petite que celle des cercles (par exemple 1,2 contre 3,0
  et 29 pour 5 sources à \(\sigma = 10^{-4}\)), car elle voit les
  sources sous toutes les directions.
\item
  \textbf{Le problème est plus dur qu'en 2D} : la source centrale est au
  milieu du cube, entourée par les sources, et une source a 4 paramètres
  au lieu de

  \begin{enumerate}
  \def\labelenumii{\arabic{enumii}.}
  \setcounter{enumii}{2}
  \tightlist
  \item
    Avec la source centrale, on retrouve de façon fiable 1 ou 2 sources
    à \(\sigma = 10^{-3}\), 2 ou 3 à \(\sigma = 10^{-4}\). Les échecs
    suivent la CRB : presque aucun tirage n'est faisable au-delà.
  \end{enumerate}
\item
  \textbf{Le BLASSO du projet échoue partout}, avec 4 à 14 sources
  estimées : noyau brut (\(\alpha = 0\)), donc le même défaut qu'en 2D
  (fausses sources du côté des capteurs). Ce n'est pas une comparaison à
  \(\lambda\) optimisé, mais elle reprend exactement le réglage du
  script du projet.
\item
  \textbf{Coût} : 0,2 à 0,7 s pour une source, 1 à 2 s pour deux, 3 s
  pour trois, 7 à 11 s pour cinq (grille de \(21^3\) nœuds, 160
  capteurs).
\end{enumerate}

\subsection{7. Limites}\label{limites}

\begin{enumerate}
\def\labelenumi{\arabic{enumi}.}
\tightlist
\item
  \textbf{Le bruit doit être connu.} La méthode utilise
  \(s_i = \sigma y_i\). Un \(\sigma\) sous-estimé d'un facteur 2 crée
  des sources parasites (section 6.4). Si \(\sigma\) est inconnu, on
  peut l'estimer par \(\chi^2/(M - pK)\) sur une première
  reconstruction, ou le surestimer légèrement.
\item
  \textbf{Modèle gaussien.} La vraisemblance suppose un bruit gaussien
  indépendant. Pour un bruit de comptage (loi de Poisson, option
  \texttt{poisson\_counts} du projet), il faudrait remplacer \(\chi^2\)
  par la déviance de Poisson ; le reste de la méthode est inchangé.
\item
  \textbf{Optimum local.} La descente conjointe n'est pas convexe. Le
  départ glouton est bon en pratique (en partant des vraies positions,
  on retrouve la même solution dans les cas vérifiés), mais rien ne
  garantit l'optimum global, surtout pour des sources proches. Les
  variantes à plusieurs candidats ou à départ depuis Prony n'ont rien
  apporté (compte rendu général, section 5.5).
\item
  \textbf{Sources faibles ou proches.} Une source est vue si son SNR
  dépasse environ 3, et localisée à 1 près si son SNR dépasse environ 20
  ; deux sources plus proches que la séparation minimale (section 6.3)
  sont fusionnées en une seule. C'est une limite des données, que la CRB
  quantifie.
\item
  \textbf{Coût.} Il croît avec \(K\) (un glissement à chaque ajout) :
  moins d'une demi-seconde en 2D jusqu'à 5 sources, mais 4,5 s avec 1024
  capteurs et 7 à 11 s en 3D pour 5 sources. Pour \(K\approx 20\)
  (configuration par défaut du projet), il faudrait limiter le
  glissement aux sources voisines de la nouvelle.
\item
  \textbf{Configuration par défaut du projet}
  (\(N\sim\mathrm{Poisson}(20)\), \(\sigma = 1\,\%\)) : la CRB conjointe
  est comprise entre \(10^4\) et \(5\times 10^5\). Aucune méthode ne
  peut retrouver ces sources ; ML + BIC en retient 1 à 5, les plus
  intenses (compte rendu général, section 5.5).
\end{enumerate}

\subsection{8. Recommandations}\label{recommandations}

{\def\LTcaptype{none} % do not increment counter
\begin{longtable}[]{@{}
  >{\raggedright\arraybackslash}p{(\linewidth - 2\tabcolsep) * \real{0.5000}}
  >{\raggedright\arraybackslash}p{(\linewidth - 2\tabcolsep) * \real{0.5000}}@{}}
\toprule\noalign{}
\begin{minipage}[b]{\linewidth}\raggedright
question
\end{minipage} & \begin{minipage}[b]{\linewidth}\raggedright
réponse
\end{minipage} \\
\midrule\noalign{}
\endhead
\bottomrule\noalign{}
\endlastfoot
quelle méthode utiliser ? & \texttt{greedy\_ml} (2D) ou
\texttt{greedy\_ml\_nd} (3D), critère BIC \\
quel bruit lui donner ? & \(s_i = \sigma y_i\) sur les mesures brutes,
source centrale comprise (\texttt{relative\_noise\_std}) ; en cas de
doute, un \(\sigma\) un peu trop grand \\
comment savoir si le résultat est fiable ? & \(\chi^2/M\) proche de
\((M - pK)/M\), et CRB de chaque source (\texttt{cramer\_rao}) petite
devant la distance entre sources \\
\(K\) est connu à l'avance ? &
\texttt{k\_max=K,\ criterion=0.0,\ backward=False} \\
comment améliorer les résultats ? & réduire le bruit de la source
centrale (c'est le facteur dominant), ajouter des capteurs (gain en
\(1/\sqrt M\)), en 3D préférer une sphère de capteurs \\
à éviter & BLASSO non pondéré (\(\alpha = 0\)), AIC, capteurs en
hauteur, un seul cercle en 3D \\
\end{longtable}
}

\subsection{9. Pistes}\label{pistes}

\begin{enumerate}
\def\labelenumi{\arabic{enumi}.}
\tightlist
\item
  \textbf{Estimer \(\sigma\) en même temps} que les sources (maximum de
  vraisemblance en \(\sigma\) :
  \(\hat\sigma^2 = \frac{1}{M}\sum_i r_i^2/y_i^2\)), pour ne plus
  dépendre d'une valeur fournie.
\item
  \textbf{Bruit de Poisson} : remplacer \(\chi^2\) par la déviance de
  Poisson (même algorithme glouton, même BIC).
\item
  \textbf{Glissement local} pour les grands \(K\) : ne réoptimiser que
  les sources proches de la dernière ajoutée, puis un glissement global
  à la fin.
\item
  \textbf{Intervalles de confiance} : la matrice \(F^{-1}\) évaluée en
  la solution donne une ellipse d'incertitude par source ; l'afficher
  dans \texttt{tester\_ml.py}.
\item
  \textbf{Plan d'expérience} : utiliser la CRB pour choisir la position
  des capteurs avant de les construire (par exemple des capteurs à
  l'intérieur de \(D\), plus proches des sources).
\end{enumerate}

\subsection{10. Reproduire}\label{reproduire}

\begin{Shaded}
\begin{Highlighting}[]
\ExtensionTok{pip}\NormalTok{ install }\AttributeTok{{-}e} \StringTok{".[dev]"}\NormalTok{ scikit{-}learn}
\ExtensionTok{pytest}\NormalTok{ sparse\_sources/tests                      }\CommentTok{\# tous les tests}
\ExtensionTok{python}\NormalTok{ sparse\_sources/experiments\_projet.py      }\CommentTok{\# comparaison des méthodes (environ 6 min)}
\ExtensionTok{python}\NormalTok{ sparse\_sources/experiments\_ml.py          }\CommentTok{\# études de ML + BIC}
\ExtensionTok{python}\NormalTok{ sparse\_sources/experiments\_cylindre.py    }\CommentTok{\# capteurs sur un cylindre}
\ExtensionTok{python}\NormalTok{ sparse\_sources/experiments\_ml3d.py        }\CommentTok{\# sources en 3D}
\ExtensionTok{python}\NormalTok{ sparse\_sources/tester\_ml.py }\AttributeTok{{-}{-}help}        \CommentTok{\# tester un tirage}
\end{Highlighting}
\end{Shaded}

Pour obtenir une version LaTeX ou PDF du rapport :

\begin{Shaded}
\begin{Highlighting}[]
\ExtensionTok{pandoc}\NormalTok{ RAPPORT\_ML\_BIC.md }\AttributeTok{{-}s} \AttributeTok{{-}o}\NormalTok{ RAPPORT\_ML\_BIC.tex}
\ExtensionTok{pandoc}\NormalTok{ RAPPORT\_ML\_BIC.md }\AttributeTok{{-}s} \AttributeTok{{-}{-}pdf{-}engine}\OperatorTok{=}\NormalTok{xelatex }\AttributeTok{{-}o}\NormalTok{ RAPPORT\_ML\_BIC.pdf}
\end{Highlighting}
\end{Shaded}

\subsection{Références}\label{ruxe9fuxe9rences}

\begin{itemize}
\tightlist
\item
  G. Schwarz, \emph{Estimating the dimension of a model}, The Annals of
  Statistics, 1978 (critère BIC).
\item
  H. Akaike, \emph{A new look at the statistical model identification},
  IEEE Transactions on Automatic Control, 1974 (critère AIC).
\item
  S. M. Kay, \emph{Fundamentals of Statistical Signal Processing, Volume
  I: Estimation Theory}, Prentice Hall, 1993 (borne de Cramér--Rao,
  efficacité du maximum de vraisemblance).
\item
  S. M. Kay, \emph{Fundamentals of Statistical Signal Processing, Volume
  II: Detection Theory}, Prentice Hall, 1998 (détection, loi du
  \(\chi^2\) décentrée).
\item
  V. Duval, G. Peyré, \emph{Exact support recovery for sparse spikes
  deconvolution}, Foundations of Computational Mathematics, 2015
  (BLASSO, condition de source non dégénérée).
\item
  Q. Denoyelle, V. Duval, G. Peyré, E. Soubies, \emph{The sliding
  Frank-Wolfe algorithm and its application to super-resolution
  microscopy}, Inverse Problems, 2019.
\end{itemize}

\end{document}
